\section{Flight Discount}
\label{sec:cses-flight-discount}

\problemstatement{A directed graph of $n$ cities and $m$ flights with
prices. You hold one discount coupon that halves the price of a single
flight. Find the minimum total price to travel from city $1$ to city $n$.}

\begin{patternbox}
``One special action, used at most once'' is a \textbf{layered Dijkstra}:
double the state to $(\text{city},\ \text{coupon used?})$. The algorithm is
unchanged; only the graph grows to two copies, with discounted edges
crossing from the ``unused'' layer to the ``used'' layer.
\end{patternbox}

\subsection*{Two copies of the graph}

Run Dijkstra over states $(u,0)$ (coupon still available) and $(u,1)$
(coupon spent). From $(u,0)$, a normal edge $u\to v$ of price $w$ leads to
$(v,0)$ at cost $w$, \emph{or} you spend the coupon to reach $(v,1)$ at cost
$w/2$. From $(u,1)$, only normal edges to $(v,1)$ are allowed. The answer is
$\textit{dist}[(n,1)]$ -- reaching $n$ having used the coupon (using it on
the cheapest beneficial edge is handled automatically by taking the
minimum).

\begin{ideabox}[Encode ``Coupon Used'' in the Node]
Rather than special-casing which edge to discount, expand the state space so
the coupon status is part of the node. Dijkstra then explores all
possibilities and the shortest path to $(n,1)$ is optimal by construction.
\end{ideabox}

\subsection*{Algorithm}

\begin{enumerate}
  \item States $(u,0)$ and $(u,1)$; start $(1,0)$ at cost $0$.
  \item Relax: $(u,0)\to(v,0)$ cost $w$; $(u,0)\to(v,1)$ cost $w/2$;
        $(u,1)\to(v,1)$ cost $w$.
  \item Answer $\textit{dist}[(n,1)]$.
\end{enumerate}

\subsection*{Dry run}

\dryrun{$1\!\to\!2(10)$, $2\!\to\!3(6)$; $n=3$.}

\begin{itemize}
  \item Discount the pricier edge $1\!\to\!2$: $5+6=11$. Discount
        $2\!\to\!3$: $10+3=13$.
  \item Minimum is $11$.
\end{itemize}

\subsection*{C++ implementation}

\begin{lstlisting}
#include <bits/stdc++.h>
using namespace std;

int main() {
    int n, m;
    cin >> n >> m;
    vector<vector<pair<int,long long>>> adj(n + 1);
    for (int i = 0; i < m; i++) {
        int a, b; long long w; cin >> a >> b >> w;
        adj[a].push_back({b, w});
    }

    // dist[u][0/1] : 0 = coupon unused, 1 = coupon used
    vector<array<long long,2>> dist(n + 1, {LLONG_MAX, LLONG_MAX});
    dist[1][0] = 0;
    priority_queue<tuple<long long,int,int>,
                   vector<tuple<long long,int,int>>, greater<>> pq;
    pq.push({0, 1, 0});
    while (!pq.empty()) {
        auto [d, u, used] = pq.top(); pq.pop();
        if (d > dist[u][used]) continue;
        for (auto& [v, w] : adj[u]) {
            if (d + w < dist[v][used]) {              // normal edge, same layer
                dist[v][used] = d + w; pq.push({dist[v][used], v, used});
            }
            if (used == 0 && d + w / 2 < dist[v][1]) { // spend coupon here
                dist[v][1] = d + w / 2; pq.push({dist[v][1], v, 1});
            }
        }
    }
    cout << dist[n][1] << "\n";
    return 0;
}
\end{lstlisting}

\begin{complexitybox}
\textbf{Time Complexity.} $O((n+m)\log n)$ -- Dijkstra on a graph of double
size. \textbf{Space Complexity.} $O(n+m)$.
\end{complexitybox}

\begin{takeawaybox}
A single one-time special move becomes an extra boolean in the state, and
plain Dijkstra on the layered graph solves it. This state-augmentation
trick generalises to ``at most $k$ specials'' (a $k+1$-layer graph) and is
one of the most reusable shortest-path patterns.
\end{takeawaybox}
